\section{The inverse problem}\label{sec:tt-question}
The experimentally meaningful question is not whether a model possesses a ``triton--triton potential.'' Rather, one asks which elastic-channel amplitudes can be reconstructed, in a stated energy interval, from measured angular distributions and other observables once Coulomb scattering, experimental response, and inelastic coupling are treated consistently.

Three levels of statements must be kept separate: an observable such as $\dd\sigma/\dd\Omega$; an amplitude or phase shift extracted with specified assumptions; and a short-range interaction model reproducing that amplitude. No single experimental distribution uniquely determines an off-shell potential.

\subsection{Scope and exclusions}
The entrance channel consists of two identical spin-$\tfrac12$, charge-$+e$ clusters. The elastic threshold, internal breakup thresholds, possible rearrangement channels and resonant structures require explicit channel bookkeeping; a resonance inferred from a different $A=6$ reaction must not be silently treated as an elastic $t+t$ measurement.

\subsection{Questions driving this review}
\begin{enumerate}[label=\arabic*.]
\item Which experiments directly measure $t+t$ elastic scattering, and over what energies and angles?
\item Which published analyses report partial-wave phase shifts or Coulomb-modified effective-range parameters, with uncertainty and normalization conventions?
\item Which six-nucleon or cluster-level Hamiltonians can predict the relevant on-shell observables, and which numerical approximations are controlled?
\item How much of an observed correlation function can be attributed to source geometry rather than the interaction amplitude?
\end{enumerate}
